\chapter*{A notation primer}
\addcontentsline{toc}{chapter}{A notation primer}
\markboth{A notation primer}{A notation primer}

Two notations run through this book: ordinary mathematics, set in the usual way,
and a \emph{pseudo-language} in which we write the simulator's algorithms. This
short chapter teaches you to read both. You do not need to memorize anything
here; treat it as a reference and come back when a symbol is unfamiliar.

\section*{Reading the pseudo-language}

The algorithms in this book are written in a small, deliberately readable
notation borrowed from typed functional programming. It is not a language you
will run; it is a language you will \emph{read}, the way a physicist reads a
Feynman diagram. It always appears in a shaded, monospaced box, like this:

\begin{lstlisting}[style=l4style]
-- a function's TYPE: "given a scalar and two tensors, return a tensor"
axpy :: Scalar -> Tensor[N] -> Tensor[N] -> Tensor[N]

-- its DEFINITION: y' = a*x + y  (computed as a fresh value, nothing overwritten)
axpy a x y = a .* x .+ y
\end{lstlisting}

A handful of conventions cover almost everything:

\begin{description}[style=nextline,leftmargin=1.5em]
  \item[\code{f :: A -> B}] is a \emph{type signature}: ``\code{f} is a function
  taking an \code{A} and producing a \code{B}.'' The arrow \code{->} reads
  ``maps to.''
  \item[Currying.] A function of several arguments is written
  \code{f :: A -> B -> C}, read ``given an \code{A}, then a \code{B}, produce a
  \code{C}.'' Arguments are supplied one at a time, separated by spaces:
  \code{f x y}. This lets us supply some arguments now and the rest later---the
  basis of ``build once, apply many times.''
  \item[Records] are bundles of named fields, written in braces:
  \lstinline[style=l4style]|{ x: Tensor[N], it: Int, converged: Bool }|. A field
  is read with a dot, \code{s.x}. A record is \emph{updated} by making a fresh
  copy with one field changed: \lstinline[style=l4style]|{ ...s, it: s.it + 1 }|
  means ``\code{s} again, but with \code{it} increased by one.'' The original
  \code{s} is untouched.
  \item[Lists] are written \code{[a, b, c]}; a \emph{list comprehension}
  \code{[ f x | x <- xs ]} reads ``the list of \code{f x} for each \code{x} in
  \code{xs}.''
  \item[\code{foldl}, \code{foldr}, \code{map}] are the standard
  \emph{combinators}: \code{map f xs} applies \code{f} to every element;
  \code{foldl f z xs} threads an accumulator left to right through the list.
  These three, and one loop combinator we introduce in \cref{ch:fold}, express
  every iteration in the book.
  \item[\code{do} blocks] sequence a computation top to bottom, like imperative
  code:
\begin{lstlisting}[style=l4style]
do  Ax  <- apply A x        -- name an intermediate result
    a   <- dot Ax x
    modify (\s -> { ...s, it: s.it + 1 })   -- the one place state changes
    return a
\end{lstlisting}
  We use \code{do} blocks only where a computation threads evolving state; the
  rest is plain function application. The full meaning of \code{do} and
  \code{modify} is built carefully in \cref{ch:monad}---for now, read a
  \code{do} block as ``these steps, in order.''
\end{description}

\begin{notationbox}
The pseudo-language is set in a \emph{box}, never inside a mathematical formula.
This is a real convention, not a stylistic one: the notation reuses symbols
(such as the \texttt{\$} sign in shape expressions) that would clash with
mathematics typesetting. When you see a shaded box, you are reading code; when
you see an inline formula like $\mat{K}\vx=\vb$, you are reading mathematics.
\end{notationbox}

\section*{Shapes, in brief}

Every tensor carries a \emph{shape}---the sizes and meanings of its axes---written
in square brackets after \code{Tensor}. We will treat shapes carefully in
\cref{ch:tensors}; for now you need only three forms:

\begin{itemize}[nosep]
  \item \code{Tensor[N]} is a flat vector of length $N$---a single column of
  degrees of freedom, the workhorse shape of every linear solve.
  \item \code{Tensor[m, n]} is a two-axis array (think: a small dense matrix).
  \item \lstinline[style=l4style]|Tensor[(S: ...)]| means ``a tensor whose axes
  form a run we will call \code{S}, of unspecified rank''; a later
  \lstinline[style=l4style]|Tensor[$S]| asserts ``the same shape \code{S}
  again.'' This lets a single function describe an operation that works on
  vectors, matrices, or higher arrays at once.
\end{itemize}

\section*{Mathematical conventions}

Vector fields are set upright and bold: $\vE$ (electric field), $\vB$ (magnetic
flux density), $\vJ$ (current density), $\vA$ (vector potential). The
differential operators are $\Grad$ (gradient), $\Div$ (divergence), and
$\Curl$ (curl). Discrete operators---matrices that act on vectors of degrees of
freedom---are set in a sans-serif font: $\mat{K}$ (a stiffness operator),
$\mat{M}$ (a mass operator), $\mat{C}$ (a damping operator). Inner products are
$\inner{\vx}{\vy}$ and norms are $\norm{\vx}$. The real and complex numbers are
$\Real$ and $\Complex$; $\Re$ and $\Im$ take real and imaginary parts. The four
finite-element function spaces, introduced in \cref{ch:derham}, are $\Hone$,
$\Hcurl$, $\Hdiv$, and $\Ltwo$.

\Cref{tab:symbols} collects the symbols used most often.

\begin{table}[t]
\centering
\caption[Frequently used symbols]{Frequently used symbols and the page or chapter
where each is introduced.}
\label{tab:symbols}
\small
\begin{tabular}{@{}lll@{}}
\toprule
Symbol & Meaning & Introduced \\
\midrule
$\vE,\vB,\vA$ & electric field, magnetic flux density, vector potential & \cref{ch:maxwell} \\
$\varepsilon,\mu,\sigma$ & permittivity, permeability, conductivity & \cref{ch:maxwell} \\
$\mat{K},\mat{M},\mat{C}$ & stiffness, mass, damping operators & \cref{ch:weak} \\
$a(u,v)$ & a bilinear (weak) form & \cref{ch:weak} \\
$\Hone,\Hcurl,\Hdiv,\Ltwo$ & the de Rham function spaces & \cref{ch:derham} \\
$N$ & number of (true) degrees of freedom & \cref{ch:fem} \\
$\omega,\lambda$ & angular frequency, eigenvalue & \cref{ch:maxwell} \\
$Q$ & quality factor of a resonant mode & \cref{ch:eigenmodes} \\
$S_{ij}$ & scattering-matrix (S-parameter) entry & \cref{ch:driven} \\
\code{Tensor[...]} & a typed, immutable array & \cref{ch:tensors} \\
\code{iterate\_while} & the loop combinator & \cref{ch:fold} \\
\code{Solve} & the state monad threading \code{SimState} & \cref{ch:monad} \\
\bottomrule
\end{tabular}
\end{table}
